\section{L1: uniform energy bounds}
\label{sec:L1}

\subsection{The continuum estimate is complete}

\begin{lemma}[Linear energy bound for the limiting norm-two field]
\label{lem:continuumEB}
In the level-one lattice VOA of the glued rank-eight lattice, the
norm-two vector $e^\lambda$ is a weight-one affine current.  For every
$f\in C^\infty(S^1)$ there is a constant $c_f$ such that
\[
 \norm{Y(e^\lambda,f)\xi}
 \le c_f\norm{(1+L_0)\xi}.
\tag{4.1}\label{eq:continuumEB}
\]
\end{lemma}

\begin{proof}
The norm formula gives $h_\lambda=\norm{\lambda}^2/2=1$.
The weight-one subspace of the lattice VOA is the direct sum of the
Cartan currents and the root vectors $e^\alpha$ for
$\norm{\alpha}^2=2$.  Their modes satisfy the level-one affine
commutation relations.  Positivity of the Sugawara expression for $L_0$
gives, for a fixed normalized current $J$,
\[
 \norm{J_n\xi}
 \le C(1+|n|)\norm{(1+L_0)^{1/2}\xi}.
\tag{4.2}\label{eq:currentmode}
\]
Indeed, on each $L_0$ eigenspace, the positive quadratic Sugawara sum
contains $J_{-n}J_n$; commuting a negative mode to the right contributes
the central term $|n|$.  Summing
$\sum_n|\widehat f_n|\norm{J_n\xi}$ and applying Cauchy--Schwarz with
$\sum_n(1+|n|)^4|\widehat f_n|^2<\infty$ gives
\[
 \norm{J(f)\xi}
 \le C\norm{f}_{H^2}\norm{(1+L_0)\xi}.
\]
This is also the weight-one case of
Theorem~\ref{thm:CKLWenergy}.
\end{proof}

\subsection{The edge-projected lattice Green form}

We now define and estimate the concrete $G_N$ used by the collar
bosonization.  It is important that this is a definition from the frozen
QWZ edge band, not a continuum ansatz.

Use antiperiodic (NS) momenta
\[
 p_n=\frac{2\pi}{N}\left(n+\frac12\right),
 \qquad -\frac N2\le n<\frac N2,
\tag{4.3}\label{eq:NSmomenta}
\]
and Hardy occupation $\nu_n=\id_{\{n<0\}}$.  The same argument works
after a fixed deck shift because only the number of indices crossing the
Hardy cut is used.  Let $u(p)$ be the normalized top-edge eigenvector of
the $2N_y$-dimensional strip matrix through $p=0$.  Away from the
crossing choose any normalized continuation; the global bound below
uses no regularity there.

For $1\le m\le N/2$, define the finite edge-current structure factor
\[
 S_N(m)
 :=\sum_{n=-m}^{-1}
 \left|\ip{u(p_{n+m})}{u(p_n)}\right|^2,
\tag{4.4}\label{eq:structurefactor}
\]
and extend it evenly to negative $m$.  With the lattice derivative
\eqref{eq:latticesymbol}, the Green multiplier is
\[
 g_N(m):=\frac{S_N(m)}{2\omega_N(m)^2}.
\tag{4.5}\label{eq:greenmultiplier}
\]
Let
$\mathbb M_N=\{-N/2+1,\ldots,-1,1,\ldots,N/2\}$ be the fundamental
nonzero Fourier representatives.  For an eight-component zero-mean
charge $q$ in the band-limited trigonometric space $\V_N$, put
\[
 \mathcal Q_N(q):=\ip{q}{G_Nq}
 =\sum_{a=1}^8\sum_{m\in\mathbb M_N}
 g_N(m)|\widehat q_a(m)|^2.
\tag{4.6}\label{eq:greenform}
\]
The continuum chiral form is
\[
 \mathcal Q_{\rm ch}(q)
 =\frac12\sum_{a=1}^8\sum_{m\ne0}
 \frac{|\widehat q_a(m)|^2}{|m|}.
\tag{4.7}\label{eq:chiralgreen}
\]
The diagonal logarithm is omitted on both sides by the same Wick
normal ordering.

\begin{lemma}[Exact transverse edge doublet]\label{lem:edgedoublet}
At $p=0$, the $M=1$, $N_y=8$ strip Hamiltonian has two zero modes, one
on each boundary.  The remaining spectrum is $\{+1,-1\}$, each with
multiplicity seven.  The two eigenbranches through zero are analytic
near $p=0$, have slopes $+1$ and $-1$, and are separated there from the
bulk complement by $\delta_{\rm bulk}=1$.
\end{lemma}

\begin{proof}
At $p=0$ the onsite term in \eqref{eq:qwz} vanishes.  The transverse
hopping is
\[
 T_y=\frac{\sigma_y}{2i}-\frac{\sigma_z}{2}.
\]
It has rank one and singular value one.  In the eigenbasis of
$\sigma_x$, the open strip is the orthogonal sum of seven two-site
blocks
$\left(\begin{smallmatrix}0&1\\1&0\end{smallmatrix}\right)$ and two
unpaired boundary vectors.  Hence the stated spectrum and
$\delta_{\rm bulk}=1$.  The derivative of the onsite matrix at zero is
$\sigma_x$; its compression to the two boundary zero modes is diagonal
with eigenvalues $+1$ and $-1$.  Analytic perturbation theory for this
isolated two-dimensional spectral subspace, followed by diagonalization
of the first derivative inside it, gives two analytic branches with
those slopes.
\end{proof}

\begin{lemma}[Sharp lattice-symbol comparison]\label{lem:sine}
For every $N$ and $0<|m|\le N/2$,
\[
 \frac{2}{\pi}|m|\le\omega_N(m)\le |m|.
\tag{4.8}\label{eq:sinecomparison}
\]
Both constants are sharp.
\end{lemma}

\begin{proof}
Set $x=\pi|m|/N\in(0,\pi/2]$.  Concavity of $\sin x$ on this interval
places its graph above the chord joining $(0,0)$ and $(\pi/2,1)$, so
$\sin x\ge2x/\pi$.  The inequality $\sin x\le x$ gives the upper bound.
The lower constant is attained at $|m|=N/2$, and the upper constant is
approached as $|m|/N\to0$.
\end{proof}

\begin{lemma}[Particle--hole Schur bound]\label{lem:Schur}
For every $N$ and $0<|m|\le N/2$,
\[
 0\le S_N(m)\le |m|.
\tag{4.9}\label{eq:Schur}
\]
\end{lemma}

\begin{proof}
For $m>0$, exactly $m$ occupied Hardy indices
$n=-m,\ldots,-1$ are carried across the Fermi cut.  Every summand in
\eqref{eq:structurefactor} lies in $[0,1]$ by Cauchy--Schwarz and the
normalization of $u(p)$.  This proves the claim; evenness handles
$m<0$.  Equivalently, this is the Schur bound for the compression of the
momentum-shift partial isometry between occupied and unoccupied edge
subspaces.
\end{proof}

\begin{lemma}[Bulk resolvent comparison]\label{lem:resolvent}
Assume the bulk-edge decomposition \eqref{eq:bulkgap}.  For
$z$ on a contour at distance $\delta_{\rm bulk}/2$ from the bulk
spectrum and for two bulk Hamiltonians $H_{\rm b},\widetilde H_{\rm b}$,
\[
 \norm{(z-H_{\rm b})^{-1}-(z-\widetilde H_{\rm b})^{-1}}
 \le\frac{4}{\delta_{\rm bulk}^2}
       \norm{H_{\rm b}-\widetilde H_{\rm b}}.
\tag{4.10}\label{eq:resolventidentity}
\]
Consequently the corresponding bulk spectral projections differ by at
most a contour-length constant times the right-hand side.
\end{lemma}

\begin{proof}
The resolvent identity gives
\[
 (z-H_{\rm b})^{-1}-(z-\widetilde H_{\rm b})^{-1}
 =(z-H_{\rm b})^{-1}(H_{\rm b}-\widetilde H_{\rm b})
   (z-\widetilde H_{\rm b})^{-1}.
\]
Both resolvent norms are at most $2/\delta_{\rm bulk}$.  Integrating the
identity in the Riesz formula proves the projection statement.
\end{proof}

\begin{remark}
Lemma~\ref{lem:resolvent} controls the massive transverse complement.
It supplies the isolated analytic edge subspace used in
Lemma~\ref{lem:edgedoublet}, but it cannot by itself furnish the
$|m|^{-1}$ bound in \eqref{eq:chiralgreen}.  That bound comes from the
exact Hardy particle--hole count and Lemma~\ref{lem:sine}.
\end{remark}

\begin{theorem}[Uniform Green-form domination and Sobolev convergence]
\label{thm:UGF}
For the edge-projected frozen $M=1$, $N_y=8$ QWZ collar, every
zero-mean $q\in\V_N\otimes\C^8$ satisfies
\[
 0\le\mathcal Q_N(q)
 \le\frac{\pi^2}{4}\mathcal Q_{\rm ch}(q).
\tag{4.11}\label{eq:UGFhere}
\]
Thus one may take the all-volume constant
\[
 K_G=\frac{\pi^2}{4}=2.467401100272\ldots .
\tag{4.12}\label{eq:KG}
\]
If $P_N$ is Fourier projection onto $|m|\le N/2$, then
$G_NP_N\to G_{\rm ch}$ in local Sobolev quadratic-form norm:
\[
 \left|\mathcal Q_N(P_Nq)-\mathcal Q_{\rm ch}(q)\right|
 \le\varepsilon_N\norm{q}_{H^2(S^1;\C^8)}^2,
 \qquad\varepsilon_N\longrightarrow0.
\tag{4.13}\label{eq:SobolevUGF}
\]
\end{theorem}

\begin{proof}
Lemmas~\ref{lem:sine} and \ref{lem:Schur} give, mode by mode,
\[
 0\le g_N(m)
 =\frac{S_N(m)}{2\omega_N(m)^2}
 \le\frac{|m|}{2(2|m|/\pi)^2}
 =\frac{\pi^2}{8|m|}
 =\frac{\pi^2}{4}\frac1{2|m|}.
\tag{4.14}\label{eq:modeUGF}
\]
Multiplication by $|\widehat q_a(m)|^2$ and summation over the modes and
eight independent components proves \eqref{eq:UGFhere}.  Notice that
the number of components does not multiply $K_G$.

It remains to prove convergence.  Fix $m$.  The $|m|$ terms in
\eqref{eq:structurefactor} involve momenta satisfying
$|p_n|+|p_{n+m}|=O_m(N^{-1})$.  By
Lemma~\ref{lem:edgedoublet}, $u(p)$ is norm-continuous there, uniformly
over this finite set.  Consequently
\[
 S_N(m)\longrightarrow |m|,
\qquad
 \frac{\omega_N(m)}{|m|}\longrightarrow1,
\qquad
 g_N(m)\longrightarrow\frac1{2|m|}.
\tag{4.15}\label{eq:pointwisegreen}
\]
Define
\[
 \varepsilon_N:=
 \sup_{0<|m|\le N/2}
 \frac{|g_N(m)-(2|m|)^{-1}|}{(1+m^2)^2}
 +\sup_{|m|>N/2}\frac{1}{2|m|(1+m^2)^2}.
\tag{4.16}\label{eq:epsilonN}
\]
For any fixed $M$, the first supremum over $|m|\le M$ tends to zero by
\eqref{eq:pointwisegreen}.  For $|m|>M$,
\eqref{eq:modeUGF} bounds the numerator by
$(K_G+1)/(2|m|)$, uniformly in $N$; hence the remaining supremum tends
to zero as $M\to\infty$.  Thus $\varepsilon_N\to0$.  Parseval's identity
and the definition of the $H^2$ norm now give
\eqref{eq:SobolevUGF}.
\end{proof}

\begin{theorem}[L1 after the proved Green estimate]
\label{thm:L1reduction}
For the edge-projected bosonization, assume only that the normalized
cocycles and zero-mode charge shifts have norm one in all four grades.
Then for every smooth $f$ the bound \eqref{eq:L1} holds with $p=1$ and
\[
 C_f\le C(\pi^2/4,\lambda)\norm{f}_{H^2},
\tag{4.17}\label{eq:Cf}
\]
independent of $N$.
\end{theorem}

\begin{proof}
Bosonic Wick ordering factors a mode of \eqref{eq:chargedfield} into the
norm-one zero-mode shift and exponentials of creation and annihilation
currents.  The normal-ordered lattice-VOA estimate is the finite-mode
version of Theorem~\ref{thm:CKLWenergy}; its constant depends on the
Gaussian contraction only through the Green quadratic form.  Theorem
\ref{thm:UGF} bounds that contraction uniformly by the continuum form
with $K_G=\pi^2/4$.  The norm-two vector has weight one, so
Lemma~\ref{lem:continuumEB} and \eqref{eq:smearedbound} give degree
$p=1$ and \eqref{eq:Cf}.  Wick scalar factors are bounded by the same
quadratic-form domination.
\end{proof}

\begin{remark}[Scope of the proved UGF]
Theorem~\ref{thm:UGF} is an all-$N$ theorem for the edge-projected
bosonization used in \eqref{eq:chargedfield}; it is not a fit to the six
tested sizes.  It does not prove that every microscopic full-CAR density
operator belongs to this bosonic edge algebra.  Excluding additional
bulk/ultraviolet local scaling operators is exactly the separate
\textup{(ALG-EXH)} condition.  A completely gapless transverse covariance
would contribute $O(N)$ particle--hole channels already at $m=1$ and
would violate \eqref{eq:UGFhere}; the isolated one-branch hypothesis is
therefore load-bearing.
\end{remark}

\paragraph{Finite audit (numerical support, not proof).}
The companion probe
\path{experiments/tfpt-discovery/mmst_ugf_bound_probe.py}
checks every displayed inequality at
$N=16,24,32,48,64,96$.  It finds
$\max_{N,m}2|m|g_N(m)=1.029399027646<K_G$, verifies
$\delta_{\rm bulk}=1$ and the unit edge slope, and rejects both a
gapless transverse mutant (whose ratio grows from $64.829$ to $384.137$)
and the false symbol constant $\omega_N(m)\ge|m|$.

\section{L2: strong convergence and local algebras}
\label{sec:L2}

\subsection{Explicit dyadic mode and Fock isometries}

Let
\[
 \mathcal K_N
 =\ell^2\!\left(\{m\in\Z:0<|m|\le N/2\}\right)\otimes\C^8
\tag{5.1}\label{eq:oneparticleboson}
\]
with the single Nyquist representative convention of
$\mathbb M_N$.  Define
\[
 j_{N,2N}(e_m\otimes e_a)=e_m\otimes e_a,
\qquad m\in\mathbb M_N,\quad 1\le a\le8.
\tag{5.2}\label{eq:modeisometry}
\]
Thus $j_{N,2N}$ is the zero-extension/Shannon embedding of the old
Fourier modes into the refined lattice.  It is an isometry and
\[
 j_{2N,4N}j_{N,2N}=j_{N,4N}.
\tag{5.3}\label{eq:modecompat}
\]
On the charge-lattice sum of symmetric bosonic Fock spaces set
\[
 V_{N,2N}
 :=\id_{\ell^2(L_{\rm charge})}\otimes\Gamma_s(j_{N,2N}).
\tag{5.4}\label{eq:Fockisometry}
\]
Second quantization preserves isometries and composition, so
$V_{N,2N}^*V_{N,2N}=1$ and
$V_{2N,4N}V_{N,2N}=V_{N,4N}$.  The charge-lattice identity makes the
same statement valid in every one of the four $\lambda$ grades.

